\documentclass[11pt]{amsart}

\usepackage{amsmath,amssymb,amsthm,mathtools}
\usepackage{geometry}
\usepackage{booktabs}
\usepackage{hyperref}
\usepackage{xcolor}
\usepackage{enumitem}
\geometry{margin=1in}
\hypersetup{colorlinks=true,linkcolor=blue!70!black,citecolor=blue!70!black,
  urlcolor=blue!70!black,
  pdftitle={A Krein-space formulation of the zero side of the Weil quadratic form on Paley-Wiener spaces},
  pdfauthor={David Alejandro Trejo Pizzo}}

\newtheorem{theorem}{Theorem}[section]
\newtheorem{proposition}[theorem]{Proposition}
\newtheorem{lemma}[theorem]{Lemma}
\newtheorem{corollary}[theorem]{Corollary}
\theoremstyle{definition}
\newtheorem{definition}[theorem]{Definition}
\newtheorem{hypothesis}[theorem]{Hypothesis}
\theoremstyle{remark}
\newtheorem{remark}[theorem]{Remark}

\newcommand{\re}{\operatorname{Re}}
\newcommand{\spec}{\operatorname{spec}}
\newcommand{\half}{\tfrac12}
\newcommand{\PW}{\mathit{PW}_d}
\newcommand{\tform}{\mathfrak t}
\newcommand{\tplus}{\mathfrak t_+}
\newcommand{\tminus}{\mathfrak t_-}
\newcommand{\ttilde}{\widetilde{\mathfrak t}_-}
\newcommand{\Hplus}{H_+}
\newcommand{\rmar}{\mathfrak a}
\newcommand{\rmq}{\mathfrak q}
\newcommand{\rmp}{\mathfrak p}
\newcommand{\Top}{\mathcal T}
\newcommand{\Kop}{K}
\newcommand{\Norm}[1]{\lVert #1 \rVert}

\title[Kre\u{\i}n-space formulation of the Weil form]{A Kre\u{\i}n-space formulation of the zero side of the Weil quadratic form on Paley--Wiener spaces}
\author{David Alejandro Trejo Pizzo}
\address{Independent Researcher}
\email{dtrejopizzo@gmail.com}


\begin{document}
\begin{abstract}
For band-limited tests $F\in\PW$ we write the zero side of the explicit formula as
$\tform(F)=\langle JEF,EF\rangle$, where $E$ is evaluation at the zeros of $\zeta$ and $J$
is the involution induced by $\rho\mapsto1-\bar\rho$; $J=I$ exactly under the Riemann
Hypothesis (RH). The positive projection is injective on the range of $E$, so this range is
the graph of an angular operator $\Kop_d$, and $\Norm{\Kop_d}\le1$ if and only if
$\tform\ge0$ on the dense test class $\mathcal D_d\subset\PW$; by Weil's criterion RH is equivalent to this for all $d$. Under a
coercivity hypothesis on the ordinates of the zeros on the critical line, which we do not
prove, $\tform$ is bounded relative to the on-line sampling norm and is represented by a
bounded self-adjoint operator. The contraction criterion is a restatement of Weil
positivity and gives no information on the sign; we make no claim of progress towards RH.
\end{abstract}

\maketitle

\tableofcontents

% ════════════════════════════════════════════════════════════════════════
\section{Introduction}\label{sec:intro}
\subsection{Weil's criterion and the sign problem}
Let $\xi(s)=\half s(s-1)\pi^{-s/2}\Gamma(s/2)\zeta(s)$ be the completed Riemann zeta
function, with nontrivial zeros $\rho$. For a suitable test function the Riemann--Weil
explicit formula equates a sum over zeros with an archimedean term, a pole term and a sum
over primes. Applied to $h=FF^{\#}$ it gives a Hermitian quadratic functional $\tform(F)$
(see \eqref{eq:EF}), and Weil's criterion \cite{Weil1952,Bombieri2000} states that
\begin{equation}\label{eq:weil}
\mathrm{RH}\iff \tform(F)\ge0\ \text{ for all } F\in\mathcal D_d \text{ and all } d>0,
\end{equation}
with $\mathcal D_d$ the test class of \S\ref{sec:setup}. Under RH one has $\tform\ge0$ on
$\mathcal D_d$ for every $d$; for a \emph{fixed} $d$ the positivity of $\tform$ on
$\mathcal D_d$ is Weil positivity restricted to tests with $\operatorname{supp}\widehat
h\subset[-2d,2d]$, which is implied by RH and is not known to imply it. For small $d$ it is
known unconditionally \cite{Yoshida}. On-line zeros contribute $|F(\gamma)|^{2}\ge0$,
while off-line zeros, if they exist, contribute terms of either sign. This paper writes
this structure in operator-theoretic form.

\subsection{Main results}
We work on band-limited tests of exponential type $d>0$. Band-limited functions are
entire, so the evaluations $F(t+ib)$ at the spectral points of off-line zeros are defined.
All objects depend on $d$, and we write $\Top_d$, $\Kop_d$, $E_d$ when the dependence
matters. The weighted norm $\Norm{\cdot}_w$ and the on-line sampling form $\tplus$ are
defined in \S\ref{sec:setup}.

\begin{theorem}[Bounded realization under coercivity; \S\ref{sec:thmA}]\label{thm:A}
Fix $d>0$ and assume \textnormal{(Coer$_d$)} \textnormal{(Hypothesis~\ref{def:H})}. Then
$|\tminus(F)|\le\ttilde(F)\le C\,\tplus(F)$ on $\mathcal D_d$ for a finite
$C=C(d,c)$; $\tminus$ extends to a bounded Hermitian form on $\Hplus$, $\tform$ is
represented by a bounded self-adjoint operator $\Top_d=I+B$ with $\Norm{B}\le C$, and
$\inf\spec(\Top_d)\ge1-C$. Moreover $\inf\spec(\Top_d)\ge0\iff\tform\ge0$ on $\mathcal D_d$.
\end{theorem}

\begin{theorem}[Kre\u{\i}n identity and angular operator; \S\S\ref{sec:krein}--\ref{sec:angular}]\label{thm:B}
For every $d$ and $F\in\mathcal D_d$, $\tform(F)=\langle JEF,EF\rangle
=\Norm{P_+EF}^{2}-\Norm{P_-EF}^{2}$, with $J=J^{*}$, $J^{2}=I$ and $P_\pm=\half(I\pm J)$.
$P_+$ is injective on $R(E)$, so $R(E)$ is the graph of a \textup{(}possibly
unbounded\textup{)} operator $\Kop_d$ with domain $P_+R(E)$, and
\[
\inf_{F\ne0}\frac{\tform(F)}{\Norm{P_+EF}^{2}}=1-\Norm{\Kop_d}^{2}.
\]
In particular $\Norm{\Kop_d}\le1\iff\tform\ge0$ on $\mathcal D_d$; and if
$\tform\ge-m\,\tplus$ on $\mathcal D_d$ for some $m\ge0$ then $\Norm{\Kop_d}^{2}\le1+m$.
\end{theorem}

\begin{proposition}[The contraction criterion; \S\ref{sec:nogo}]\label{thm:C}
For every $d$ and $F\in\mathcal D_d$ with $F\ne0$,
$1-\Norm{P_-EF}^{2}/\Norm{P_+EF}^{2}=\tform(F)/\Norm{P_+EF}^{2}$; hence
$\Norm{\Kop_d}\le1$ if and only if $\rmar+\rmq-\rmp\ge0$ on $\mathcal D_d$. Consequently
$\mathrm{RH}\iff\Norm{\Kop_d}\le1$ for every $d>0$.
\end{proposition}

\subsection{Relation to prior work and scope}
The criterion ``$\Norm{\Kop_d}\le1$'' has the shape of the eigenvalue-below-$1$ device of
Connes--Consani \cite{ConnesConsani}, who obtain Weil positivity at the archimedean place
for test functions of small support from a compression of the scaling action onto Sonin's
space. The completion of the Weil form into a reproducing-kernel space is the theme of
Suzuki \cite{Suzuki}, who obtains a de~Branges space under RH. Our construction does not
assume RH, so the metric on the zero side is indefinite and the natural language is that of
Kre\u{\i}n spaces. We do not claim an identification of $\Kop_d$ with the Connes--Consani
operator; $\Kop_d$ is defined through the off-line zeros and cannot be computed without
knowing them.

Theorem~\ref{thm:A} is conditional on (Coer$_d$), which we do not prove
(\S\ref{sec:Hstatus}). Theorem~\ref{thm:B} and Proposition~\ref{thm:C} are unconditional;
Proposition~\ref{thm:C} is a restatement of Weil positivity at bandwidth $d$. Nothing in
the paper bears on the sign of $\tform$.

% ════════════════════════════════════════════════════════════════════════
\section{Setup and notation}\label{sec:setup}

\paragraph{Test class.} For $d>0$ let
$\PW=\{F\in L^{2}(\mathbb R): \widehat F=0 \text{ a.e.\ outside } [-d,d]\}$ be the
Paley--Wiener space (entire functions of exponential type $\le d$ in $L^{2}$), with
$\widehat F(\xi)=\int F(t)e^{-it\xi}\,dt$, so that
$\Norm{F}_2^{2}=\tfrac1{2\pi}\int|\widehat F|^{2}$. We use the dense subspace
\[
\mathcal D=\mathcal D_d:=\{F\in\PW:\ \widehat F\in C_c^\infty(-d,d)\};
\]
its elements and all their derivatives are $O(|t|^{-N})$ on $\mathbb R$ for every $N$, and
$\mathcal D_d\subset\mathcal D_{d'}$ for $d<d'$. By Plancherel,
\begin{equation}\label{eq:bernstein}
\text{(Bernstein)}\quad \Norm{F'}_{2}\le d\,\Norm{F}_{2}\qquad(F\in\PW).
\end{equation}
We also use the weighted norm
\[
\Norm{F}_{w}^{2}:=\int_{\mathbb R}|F(t)|^{2}\log(2+|t|)\,dt,
\]
finite on $\mathcal D$.

\paragraph{Zeros, depths, partners.} Write a nontrivial zero as
$\rho=\half+b_\rho+it_\rho$, $b_\rho,t_\rho\in\mathbb R$, $|b_\rho|<\half$, so
$b_\rho=0\iff\rho$ lies on the critical line. Zeros are counted with multiplicity
throughout. The functional equation gives the quartet
$\{\rho,\bar\rho,1-\rho,1-\bar\rho\}$. Define the partner involution
\begin{equation}\label{eq:sigma}
\sigma(\rho):=1-\bar\rho=\half-b_\rho+it_\rho,\qquad
\sigma^{2}=\mathrm{id},\quad b_{\sigma(\rho)}=-b_\rho,\ t_{\sigma(\rho)}=t_\rho.
\end{equation}
Let $N(T)$ be the number of zeros with $0<t_\rho\le T$,
$\bar N(T)=\tfrac T{2\pi}\log\tfrac T{2\pi}-\tfrac T{2\pi}+\tfrac78$ and $S(T):=N(T)-\bar
N(T)$; then $S(T)=O(\log T)$ \cite[Thms.~9.3, 9.4]{Titchmarsh}, so there is $A$ with
$|S(t)|\le A\log(2+|t|)$ for $t>0$, and the number of zeros with $t_\rho\in[u,u+1]$ is
$\ll\log(2+|u|)$. We write $\varrho(t):=\tfrac1{2\pi}\log\tfrac{|t|}{2\pi}$ for the
density ($\varrho=\bar N'$ for $t>0$). An interval of length $\pi/d$ at height $T$ contains
$O(\log T)$ zeros, and $(1+o(1))\varrho(T)\pi/d$ zeros under RH.

\paragraph{Weil functional.} For $F\in\mathcal D$ put $F^{\#}(z)=\overline{F(\bar z)}$,
$h=FF^{\#}$ and $g(x)=\tfrac1{2\pi}\int h(r)e^{-irx}\,dr$, which is smooth and supported
in $[-2d,2d]$. The explicit formula \cite{Weil1952}, \cite{Bombieri2000},
\cite[Thm.~5.12]{IwaniecKowalski} gives
\begin{equation}\label{eq:EF}
\tform(F):=\sum_\rho h(t_\rho+ib_\rho)
=\rmq(F)+\rmar(F)-\rmp(F),
\end{equation}
where $\rmq(F)=h(\tfrac i2)+h(-\tfrac i2)$,
\[
\rmar(F)=\tfrac1{2\pi}\int h(r)\,\re\tfrac{\Gamma'}{\Gamma}\big(\tfrac14+\tfrac{ir}2\big)\,dr-g(0)\log\pi,
\qquad
\rmp(F)=\sum_{n\le e^{2d}}\tfrac{\Lambda(n)}{\sqrt n}\,\big(g(\log n)+g(-\log n)\big);
\]
only prime powers $n\le e^{2d}$ occur. The spectral point of $\rho$ is $-i(\rho-\half)=
t_\rho-ib_\rho$; since the zero set is invariant under $\sigma$, summing $h$ over the points
$t_\rho+ib_\rho$ gives the same value. The sum converges absolutely by
Lemma~\ref{lem:D} and Cauchy--Schwarz, because $|h(t+ib)|=|F(t+ib)||F(t-ib)|$. Grouping the
zeros into partner pairs gives
\begin{equation}\label{eq:onoff}
\tform(F)=\underbrace{\sum_{b_\rho=0}|F(t_\rho)|^{2}}_{\tplus(F)\,\ge0}
\;+\;\underbrace{\sum_{\text{off-line pairs}}2\re\!\big[F(t_\rho+ib_\rho)\overline{F(t_\rho-ib_\rho)}\big]}_{\tminus(F)} .
\end{equation}
With $\ttilde(F):=\sum_{b_\rho\ne0}|F(t_\rho+ib_\rho)|^{2}$, the inequality
$2\re(F_+\overline{F_-})\le|F_+|^{2}+|F_-|^{2}$ summed over pairs gives
$|\tminus(F)|\le\ttilde(F)$.

\paragraph{Realization space.} $\tplus$ is a norm on $\mathcal D$ by
Lemma~\ref{lem:inj}. Let $\Hplus$ be the completion of $\mathcal D$ in the norm
$\Norm{F}_{\Hplus}^{2}:=\tplus(F)$ (the on-line sampling space).

% ════════════════════════════════════════════════════════════════════════
\section{Four lemmas}\label{sec:lemmas}

\begin{lemma}[Integration by parts]\label{lem:A}
Let $F\in\mathcal D$, $G=|F|^{2}$, and extend $S$ to $t<0$ by $S(-t)=-S(t)$. Then,
summing over all nontrivial zeros with multiplicity,
\[
\sum_{\rho}G(t_\rho)=\int_{\mathbb R}G\,\varrho\,dt+\tfrac74\,G(0)-\int_{\mathbb R}G'(t)S(t)\,dt .
\]
\end{lemma}
\begin{proof}
On $(0,\infty)$, $dN=\varrho\,dt+dS$ with $N(0^{+})=0$, $\bar N(0^{+})=\tfrac78$, so
$S(0^{+})=-\tfrac78$; $G$ decays rapidly and $S$ grows logarithmically, so
$G(t)S(t)\to0$ as $t\to\infty$, and $\varrho$ is integrable at $0$. Stieltjes integration by
parts gives
$\sum_{t_\rho>0}G(t_\rho)=\int_0^\infty G\varrho\,dt+\tfrac78G(0)-\int_0^\infty G'S\,dt$.
The zero set is symmetric under $t_\rho\mapsto-t_\rho$, there is no zero with $t_\rho=0$,
and applying the same identity to $G(-\cdot)$ with $S(-t)=-S(t)$ gives the contribution of
$t_\rho<0$.
\end{proof}

\begin{lemma}[Bound for the error term]\label{lem:B}
For $F\in\mathcal D$ and $G=|F|^{2}$,
\[
\Big|\int G'S\,dt\Big|\le 2A\int|F'(t)||F(t)|\log(2+|t|)\,dt\le 2A\,\Norm{F'}_w\Norm{F}_w .
\]
\end{lemma}
\begin{proof}
$|G'|\le2|F'||F|$ and $|S(t)|\le A\log(2+|t|)$ (also for $t<0$ by oddness); then
Cauchy--Schwarz.
\end{proof}

\begin{remark}\label{rem:Onoto}
Lemmas~\ref{lem:A}--\ref{lem:B} compare the sampling sum over all zeros with the weighted
integral $\int G\varrho$ up to an error of size $O(\Norm{F'}_w\Norm{F}_w)$. A comparison
with $\tplus$ itself, which samples only the on-line zeros, requires a coercivity bound of
the type of Hypothesis~\ref{def:H}.
\end{remark}

\begin{lemma}[Vertical $L^{2}$ bound]\label{lem:C}
For $F\in\PW$ and $|b|\le\half$,
$\Norm{F(\cdot+ib)}_2^{2}=\tfrac1{2\pi}\int_{-d}^{d}|\widehat F(\xi)|^{2}e^{-2b\xi}d\xi\le
e^{2d|b|}\Norm{F}_2^{2}\le e^{d}\Norm{F}_2^{2}.$
\end{lemma}
\begin{proof}
$F(t+ib)=\tfrac1{2\pi}\int\widehat F(\xi)e^{-b\xi}e^{it\xi}d\xi$; Parseval, then
$e^{-2b\xi}\le e^{2d|b|}$ on $|\xi|\le d$.
\end{proof}
This is an upper bound; it gives no lower bound for any quantity attached to the zeros.

\begin{lemma}[Sum over the zeros]\label{lem:D}
There is $C_d$ such that for all $F\in\mathcal D$,
$\sum_{\rho}|F(t_\rho+ib_\rho)|^{2}\le C_d\Norm{F}_w^{2}$. In particular
$\ttilde(F)\le C_d\Norm{F}_w^{2}$ and $\tplus(F)\le C_d\Norm{F}_w^{2}$.
\end{lemma}
\begin{proof}
Fix $\chi\in C_c^\infty(\mathbb R)$ with $\chi=1$ on $[-d,d]$ and let $k_b$ have Fourier
transform $\chi(\xi)e^{-b\xi}$. Then $F(\cdot+ib)=F*k_b$, and since the functions
$\chi(\xi)e^{-b\xi}$, $|b|\le\half$, are bounded in $C^{3}$ with common compact support,
$|k_b(u)|\le A_d(1+|u|)^{-3}$ uniformly for $|b|\le\half$. By Cauchy--Schwarz
$|F(t+ib)|^{2}\le\Norm{k_b}_1\int|F(v)|^{2}|k_b(t-v)|\,dv$. Since each unit interval
$[u,u+1]$ contains $\ll\log(2+|u|)$ zeros,
$\sum_\rho|k_{b_\rho}(t_\rho-v)|\ll_d\sum_{j\ge0}\log(2+|v|+j)(1+j)^{-3}\ll\log(2+|v|)$,
and the claim follows by summing over $\rho$ and integrating in $v$.
\end{proof}

\begin{remark}
The weight $\log(2+|t|)$ cannot be removed from the corresponding bound for the
real ordinates. Choose $F\in\mathcal D$, $F\ne0$, with
$\Norm{F'}_2\le\varepsilon\Norm{F}_2$, $\varepsilon=1/(8\pi A)$ (take $\widehat F$ supported
in $[-\varepsilon',\varepsilon']$, $\varepsilon'=\min(\varepsilon,d/2)$, and use
\eqref{eq:bernstein} with $\varepsilon'$ in place of $d$), and put $F_T=F(\cdot-T)$. Then
$\Norm{F_T}_2$ is constant, while by Lemmas~\ref{lem:A}--\ref{lem:B}
$\sum_\rho|F_T(t_\rho)|^{2}\ge\big(\tfrac1{2\pi}-2A\varepsilon+o(1)\big)\log T\,\Norm{F}_2^{2}
\ge\tfrac1{8\pi}\log T\,\Norm{F}_2^{2}$ for large $T$. Thus the sampling map
$F\mapsto(F(t_\rho))_\rho$ over the ordinates of all zeros, which is bounded from
$(\mathcal D,\Norm{\cdot}_w)$ into $\ell^{2}$ by the proof of Lemma~\ref{lem:D} with
$b=0$, is not bounded from $(\mathcal D,\Norm{\cdot}_2)$; under RH this map is $E$. The bound of Lemma~\ref{lem:D} counts all zeros in a window and does not use
the sparsity of off-line zeros, which may cluster locally.
\end{remark}

% ════════════════════════════════════════════════════════════════════════
\section{Coercivity}\label{sec:H}

A lower bound $\tplus(F)\ge c\Norm{F}_w^{2}$ is a weighted lower sampling inequality for the
on-line ordinates. By Beurling's density theorem \cite{Beurling,Seip} a separated real
sequence is sampling for $\PW$ when its lower uniform density exceeds $d/\pi$; the
unweighted inequality $\tplus(F)\ge c\Norm{F}_2^{2}$ would therefore follow if the on-line
ordinates contained a separated subsequence of lower uniform density $>d/\pi$, which is not
known unconditionally. The weighted inequality needs in addition a number of well-spread
on-line ordinates proportional to $\log T$ in every window at height $T$.

\begin{hypothesis}[Coercivity, (Coer$_d$)]\label{def:H}
There is $c>0$ with $\tplus(F)\ge c\,\Norm{F}_w^{2}$ for all $F\in\mathcal D_d$.
\end{hypothesis}

\begin{remark}[A counting condition]\label{lem:E}
A natural sufficient condition, whose implication to (Coer$_d$) we do not prove, is local
density of on-line zeros: there are $\eta>0$ and $T_0$ such that every interval $I$ of
length $\pi/d$ centred at a height $T\ge T_0$ satisfies
$\#\{\rho:\ b_\rho=0,\ t_\rho\in I\}\ge\eta\,\varrho(T)\,|I|$. Deriving (Coer$_d$) from it
would require a weighted sampling inequality that controls clustering of these ordinates
inside each window, and a separate treatment of heights $|t|<T_0$.
\end{remark}

\subsection{Remarks on the hypothesis}\label{sec:Hstatus}
Under RH, $S(T)=O(\log T/\log\log T)$ (Littlewood; see \cite[Thm.~14.13]{Titchmarsh}), so
every interval of length $\pi/d$ at height $T$ contains $(1+o(1))\varrho(T)\pi/d$ zeros,
all on the line; thus the counting condition of Remark~\ref{lem:E} holds under RH.
Unconditionally it is open: it requires both a lower bound $\gg\log T$ for the number of
zeros in every interval of fixed length (not known, since $S(T+h)-S(T)=O(\log T)$ only) and
a short-interval bound for zeros off the line. Selberg's density estimate
$N(\sigma,T)\ll T^{1-(\sigma-1/2)/4}\log T$ \cite{Selberg1946} does not exclude
$\gg\log T$ zeros with $b_\rho\ge\delta$ in one interval of length $\pi/d$.
Heuristically (extrapolating Selberg's central limit theorem for $S$
\cite{Selberg,Selberg1946}) such intervals are rare.

Without a coercivity bound we cannot exclude sequences $F_n\in\mathcal D$ with
$\tplus(F_n)\to0$ and $\tminus(F_n)\not\to0$, in which case $\tminus$ does not extend
continuously to $\Hplus$. We do not know whether (Coer$_d$) is necessary for this. In the
norm $\Norm{EF}_{\ell^2}$ of all zeros the form is bounded unconditionally,
$|\tform(F)|\le\Norm{EF}^{2}$, and the question becomes the sign alone.

% ════════════════════════════════════════════════════════════════════════
\section{Proof of Theorem~\ref{thm:A}}\label{sec:thmA}

\begin{proof}[Proof of Theorem~\ref{thm:A}]
By \S\ref{sec:setup}, $|\tminus|\le\ttilde$. By Lemma~\ref{lem:D} and (Coer$_d$),
$\ttilde\le C_d\Norm{F}_w^{2}\le (C_d/c)\,\tplus=:C\,\tplus$. Hence $\tminus$ is a
Hermitian form on $\mathcal D$ with $|\tminus(F)|\le C\Norm{F}_{\Hplus}^{2}$; by
polarization it is bounded as a sesquilinear form, and since $\mathcal D$ is dense in
$\Hplus$ it extends uniquely to a bounded Hermitian form on $\Hplus$. By the Riesz
representation theorem there is a bounded self-adjoint $B$ on $\Hplus$ with
$\tminus(F)=\langle BF,F\rangle_{\Hplus}$ and $\Norm{B}\le C$. Then
$\tform(F)=\langle(I+B)F,F\rangle_{\Hplus}$ on $\mathcal D$, so $\Top_d:=I+B$ represents
$\tform$ and $\inf\spec(\Top_d)\ge1-C$. Finally $\inf\spec(\Top_d)=\inf_{F\ne0}
\langle\Top_dF,F\rangle/\Norm{F}^{2}_{\Hplus}$, and by density and continuity this infimum
may be taken over $F\in\mathcal D$, which gives the last assertion.
\end{proof}

\begin{corollary}[Reformulation]\label{cor:Aprime}
The positivity $\tform\ge0$ on $\mathcal D_d$ does not refer to any inner product. Hence
RH implies $\inf\spec(\Top_d)\ge0$ for every $d$ for which (Coer$_d$) holds; and if
(Coer$_d$) holds for every $d>0$, then
\[
\mathrm{RH}\iff\inf\spec(\Top_d)\ge0\ \text{ for every } d>0 .
\]
\end{corollary}
\begin{proof}
Theorem~\ref{thm:A} and \eqref{eq:weil}.
\end{proof}

The constant $C$ produced by Lemma~\ref{lem:D} and (Coer$_d$) is an upper bound obtained
without using cancellation; the method gives no value of $C$ below $1$, so the resulting
lower bound $1-C$ for the bottom is not known to be non-negative. Theorem~\ref{thm:A} does
not determine the sign of the bottom.

% ════════════════════════════════════════════════════════════════════════
\section{The Kre\u{\i}n identity}\label{sec:krein}

Define the evaluation operator $E:\mathcal D\to\ell^{2}(\rho)$,
$(EF)_\rho:=F(t_\rho+ib_\rho)$, which is bounded for $\Norm{\cdot}_w$ by
Lemma~\ref{lem:D}, and the linear involution $J$ induced by the partner map
\eqref{eq:sigma}, $(Jc)_\rho:=c_{\sigma(\rho)}$. Then $J^{2}=I$ and $J=J^{*}$, since $J$ is
the permutation of coordinates by an involution; $J$ is a fundamental symmetry and
$[a,b]_J:=\langle Ja,b\rangle$ is a Kre\u{\i}n-space inner product on $\ell^{2}(\rho)$.

\begin{proposition}[The identity]\label{prop:identity}
With $\langle a,b\rangle=\sum_\rho a_\rho\overline{b_\rho}$,
\[
\langle EF,JEF\rangle=\sum_\rho (EF)_\rho\overline{(EF)_{\sigma(\rho)}}
=\sum_\rho F(t_\rho+ib_\rho)\overline{F(t_\rho-ib_\rho)}=\tform(F).
\]
\end{proposition}
\begin{proof}
Immediate from $b_{\sigma(\rho)}=-b_\rho$, $t_{\sigma(\rho)}=t_\rho$ and \eqref{eq:EF};
$\tform(F)$ is real, so $\langle EF,JEF\rangle=\langle JEF,EF\rangle$.
\end{proof}

\begin{remark}
For real-type $F$ ($F(\bar z)=\overline{F(z)}$) one has
$(JEF)_\rho=F(t_\rho-ib_\rho)=\overline{(EF)_\rho}$, so on such vectors $J$ acts as
complex conjugation.
\end{remark}

\paragraph{Three structural consequences.}
\begin{enumerate}[label=\textbf{(C\arabic*)}]
\item \emph{RH $\Rightarrow J=I$.} Under RH all $b_\rho=0$, so $\sigma(\rho)=\rho$,
$J=I$ and $\tform(F)=\Norm{EF}^{2}\ge0$.
\item \emph{If RH fails,
$J=I_{\mathrm{on}}\oplus\bigoplus_{\text{pairs}}\left(\begin{smallmatrix}0&1\\1&0\end{smallmatrix}\right)$.}
Each off-line quartet splits under $\sigma$ into two pairs $\{\rho,1-\bar\rho\}$ and
$\{\bar\rho,1-\rho\}$, on each of which $J$ is a $2\times2$ swap with eigenvalues $\pm1$. The
negative eigenspace of $J$ is spanned by vectors supported on off-line zeros.
\item \emph{RH $\Rightarrow$ $R(E_d)$ is $J$-non-negative for every $d$; conversely, if
$R(E_d)$ is $J$-non-negative for every $d>0$, then RH} (by \eqref{eq:weil}).
\end{enumerate}

% ════════════════════════════════════════════════════════════════════════
\section{The angular operator}\label{sec:angular}

Let $\ell^{2}(\rho)=\mathcal K_+\oplus\mathcal K_-$ with $J=P_+-P_-$. On a pair the $\pm1$
eigenvectors are $\tfrac1{\sqrt2}(e_\rho\pm e_{\sigma\rho})$, so with
$F_\pm:=F(t_\rho\pm ib_\rho)$,
\[
\Norm{P_+EF}^{2}=\tplus+\!\!\sum_{\text{pairs}}\tfrac12|F_++F_-|^{2},\qquad
\Norm{P_-EF}^{2}=\!\!\sum_{\text{pairs}}\tfrac12|F_+-F_-|^{2},
\]
and
\begin{equation}\label{eq:split}
\tform(F)=\Norm{P_+EF}^{2}-\Norm{P_-EF}^{2}.
\end{equation}

\begin{lemma}[Injectivity]\label{lem:inj}
If $F\in\PW$ vanishes at all ordinates $t_\rho$ of zeros on the critical line, then $F=0$.
In particular $P_+EF=0\Rightarrow F=0$ for $F\in\mathcal D$.
\end{lemma}
\begin{proof}
By Selberg's theorem \cite{Selberg1942} the number $N_0(T)$ of distinct zeros
$\half+i\gamma$, $0<\gamma\le T$, satisfies $N_0(T)\gg T\log T$ (Selberg's method counts
sign changes of Hardy's function), and the on-line ordinates are symmetric
under $\gamma\mapsto-\gamma$. Suppose $F\ne0$; dividing by a power of $z$ we may assume
$F(0)\ne0$, which preserves exponential type $d$ and the vanishing at $\pm\gamma$ (since
$\gamma\ne0$). As $|F(z)|\le C e^{d|z|}$, Jensen's formula gives
$\int_0^{er}n_F(s)\,\tfrac{ds}{s}\le d\,er+O(1)$, hence the number $n_F(r)$ of zeros of
$F$ in $|z|\le r$ satisfies $n_F(r)\le e\,d\,r+O(1)$. This contradicts $n_F(r)\ge 2N_0(r)\gg
r\log r$. The on-line coordinates of $P_+EF$ are the values $F(t_\rho)$, which gives the
second statement.
\end{proof}

Hence $R(E)$ is the graph of the \emph{angular operator}
$\Kop=\Kop_d:=P_-(P_+|_{R(E)})^{-1}$ with domain $P_+R(E)\subset\mathcal K_+$ and values in
$\mathcal K_-$ (see \cite[Ch.~1, \S8]{AzizovIokhvidov} for angular operators). We do not
assert that $\Kop_d$ is closable.

\begin{proof}[Proof of Theorem~\ref{thm:B}]
The identity is Proposition~\ref{prop:identity} and \eqref{eq:split}; injectivity is
Lemma~\ref{lem:inj}. For $x=P_+EF\ne0$ we have $\Kop x=P_-EF$, so
$\tform(F)/\Norm{P_+EF}^{2}=1-\Norm{\Kop x}^{2}/\Norm{x}^{2}$, and taking the infimum gives
$1-\Norm{\Kop_d}^{2}$. If $\tform\ge-m\tplus$ with $m\ge0$, then since $\tplus\le\Norm{P_+EF}^{2}$,
$\Norm{P_-EF}^{2}=\Norm{P_+EF}^{2}-\tform(F)\le(1+m)\Norm{P_+EF}^{2}$.
\end{proof}

In summary,
\begin{equation}\label{eq:dict}
\begin{gathered}
\Kop_d\ \text{is defined (Lemma~\ref{lem:inj})};\qquad
\tform\ge-m\,\tplus\ (m\ge0)\ \Rightarrow\ \Norm{\Kop_d}^{2}\le1+m;\\
\Norm{\Kop_d}\le1\iff\tform\ge0\ \text{on}\ \mathcal D_d .
\end{gathered}
\end{equation}
The converse of the middle implication holds under (Coer$_d$): if
$\Norm{\Kop_d}<\infty$, then $\tform\ge(1-\Norm{\Kop_d}^{2})\Norm{P_+EF}^{2}$, and
$\Norm{P_+EF}^{2}\le\tplus+\ttilde\le(1+C)\tplus$ by Lemma~\ref{lem:D} and (Coer$_d$).
Under (Coer$_d$), Lemma~\ref{lem:D} gives $\Norm{\Kop_d}^{2}\le C$, an upper bound that
is not shown to be $\le1$; no lower bound for $\Norm{\Kop_d}$ is known (under RH,
$\Kop_d=0$).

\subsection{The negative coordinates}
On a pair the negative coordinate is the vertical difference
$\Delta_\rho:=F(t_\rho+ib_\rho)-F(t_\rho-ib_\rho)=-2\,(\sinh(b_\rho D)F)(t_\rho)$ with
$D=-i\,d/dt$; that is, $F(\cdot+ib)-F(\cdot-ib)$ has Fourier multiplier $-2\sinh(b\xi)$.
Since $1\le\sinh(b\xi)/(b\xi)\le\sinh(d/2)/(d/2)$ for $|b|\le\half$, $|\xi|\le d$, the
function $F(\cdot+ib)-F(\cdot-ib)$ is comparable in $L^{2}$ to $2ibF'$ with constants
depending only on $d$. The leading term is $2ib_\rho F'(t_\rho)$, and since
$0\le\sinh x-x\le\tfrac{x^{3}}6\cosh x$, the pointwise error is at most
$\tfrac{|b_\rho|^{3}}{3}\cosh(\tfrac d2)\cdot\tfrac1{2\pi}\int|\xi|^{3}|\widehat F(\xi)|\,d\xi$.

\begin{remark}
In these coordinates
\begin{equation}\label{eq:rayleigh}
\Norm{\Kop_d}^{2}=\sup_{F\ne0}\frac{\sum_{\text{pairs}}\tfrac12|F(t_\rho+ib_\rho)-F(t_\rho-ib_\rho)|^{2}}
{\tplus(F)+\sum_{\text{pairs}}\tfrac12|F(t_\rho+ib_\rho)+F(t_\rho-ib_\rho)|^{2}},
\end{equation}
and $\mathrm{RH}\iff\Norm{\Kop_d}\le1$ for every $d>0$.
\end{remark}

% ════════════════════════════════════════════════════════════════════════
\section{The contraction criterion is a restatement of Weil positivity}\label{sec:nogo}

\begin{proof}[Proof of Proposition~\ref{thm:C}]
By \eqref{eq:split} and \eqref{eq:EF},
\begin{equation}\label{eq:pivot}
\Norm{P_+EF}^{2}-\Norm{P_-EF}^{2}=\tform(F)=\rmar(F)+\rmq(F)-\rmp(F).
\end{equation}
Dividing by $\Norm{P_+EF}^{2}>0$ (Lemma~\ref{lem:inj}) gives the first identity, and
Theorem~\ref{thm:B} gives $\Norm{\Kop_d}\le1\iff\tform\ge0\iff\rmar+\rmq-\rmp\ge0$ on
$\mathcal D_d$. The last statement follows from \eqref{eq:weil}.
\end{proof}

\begin{remark}\label{cor:circular}
The criterion $\Norm{\Kop_d}\le1$ is therefore a restatement of Weil positivity at
bandwidth $d$ and contains no additional information. Only the condition
$\Norm{\Kop_d}\le1$, not the value of $\Norm{\Kop_d}$, is determined by the arithmetic side:
$\Norm{P_+EF}^{2}$ depends on the depths $b_\rho$ separately. Moreover $\Kop_d$ is defined
through the off-line zeros, so it cannot be computed without knowing them. Bounds for
$\Norm{P_-EF}$ that ignore the relative phase of $F(t_\rho\pm ib_\rho)$, such as
Lemma~\ref{lem:D}, give under (Coer$_d$) upper bounds for $\Norm{\Kop_d}$ that are not
shown to be $\le1$.
\end{remark}

% ════════════════════════════════════════════════════════════════════════
\section{Relation to prior work}\label{sec:prior}

\paragraph{Connes--Consani \cite{ConnesConsani}.} They relate the archimedean part of the
Weil functional to the compression of the scaling action on $L^{2}(\mathbb
R)_{\mathrm{ev}}$ onto Sonin's space, and prove positivity of the archimedean Weil functional
for test functions supported in $[2^{-1/2},2^{1/2}]$ whose Fourier transform satisfies the
vanishing conditions of \cite[Thm.~1.1]{ConnesConsani}, using prolate spheroidal wave functions and an
eigenvalue-below-$1$ argument. Our criterion $\Norm{\Kop_d}\le1$ has a similar shape. We
do not claim an identification of $\Kop_d$ with their operator: the two are distinct
objects, ours defined through the zeros, theirs through archimedean data independent of
the zeros.

\paragraph{Suzuki \cite{Suzuki}.} Completing the Weil Hermitian form yields, under RH, a
de~Branges space attached to $\xi$ (see also \cite{Lagarias} for de~Branges structures for
$L$-functions); Suzuki derives from it a condition equivalent to RH. On the zero side,
with off-line zeros allowed, the metric of \S\ref{sec:krein} is indefinite, so the natural
object is a Kre\u{\i}n space (a Pontryagin space if the off-line set were finite, which is
not known).

\paragraph{Classical background.} Weil's criterion \cite{Weil1952,Bombieri2000};
de~Branges spaces \cite{deBranges}; angular operators in Kre\u{\i}n spaces
\cite{AzizovIokhvidov}; sampling in Paley--Wiener spaces \cite{PlancherelPolya,Beurling,Seip};
unconditional Weil positivity for tests of small support \cite{Yoshida}. The new content of
this paper is the rewriting $\tform=\langle JEF,EF\rangle$ and the angular-operator
restatement.

% ════════════════════════════════════════════════════════════════════════
\section{Summary and open problems}\label{sec:open}

\begin{center}
\small
\begin{tabular}{ll}
\toprule
Statement & Status \\
\midrule
Lemmas~\ref{lem:A}--\ref{lem:D} (weighted form) & proved \\
Counting condition $\Rightarrow$ (Coer$_d$) (Remark~\ref{lem:E}) & not proved \\
Theorem~\ref{thm:A} (bounded $\Top_d$, $\inf\spec\ge1-C$) & conditional on (Coer$_d$) \\
Corollary~\ref{cor:Aprime} & conditional on (Coer$_d$) for all $d$ \\
Proposition~\ref{prop:identity}: $\tform=\langle JEF,EF\rangle$ & proved \\
Theorem~\ref{thm:B}: angular operator, \eqref{eq:dict} & proved \\
Proposition~\ref{thm:C}: restatement of Weil positivity & proved \\
Sign of $\tform$ & not addressed \\
\bottomrule
\end{tabular}
\end{center}

\paragraph{Open problems.}
\begin{enumerate}[label=\textbf{\arabic*.}]
\item Does (Coer$_d$) hold unconditionally for some $d$? Does the counting condition of
Remark~\ref{lem:E} imply it?
\item Sharp constants in Lemma~\ref{lem:D}, and a quantitative weighted sampling inequality
for the on-line ordinates.
\item The sign of $\tform$ on $\mathcal D_d$ for large $d$. The zero-side formulation
studied here does not address it; a factorization $\Top_d=A^{*}A$ or a positivity
statement on the arithmetic side would be needed.
\end{enumerate}

% ════════════════════════════════════════════════════════════════════════
\begin{thebibliography}{99}

\bibitem{AzizovIokhvidov} T.~Ya.~Azizov, I.~S.~Iokhvidov, \emph{Linear Operators in
Spaces with an Indefinite Metric}, Wiley, Chichester, 1989.

\bibitem{Beurling} A.~Beurling, \emph{The Collected Works of Arne Beurling}, Vol.~2,
Birkh\"auser, Boston, 1989.

\bibitem{Bombieri2000} E.~Bombieri, \emph{Remarks on Weil's quadratic functional in the
theory of prime numbers I}, Rend.\ Mat.\ Acc.\ Lincei (9) 11 (2000), 183--233.

\bibitem{ConnesConsani} A.~Connes, C.~Consani, \emph{Weil positivity and trace formula,
the archimedean place}, Selecta Math.\ (N.S.) 27 (2021), no.~4, Paper No.~77.

\bibitem{deBranges} L.~de~Branges, \emph{Hilbert Spaces of Entire Functions},
Prentice-Hall, Englewood Cliffs, 1968.

\bibitem{IwaniecKowalski} H.~Iwaniec, E.~Kowalski, \emph{Analytic Number Theory}, Amer.\
Math.\ Soc.\ Colloq.\ Publ.\ 53, 2004.

\bibitem{Lagarias} J.~C.~Lagarias, \emph{Hilbert spaces of entire functions and Dirichlet
$L$-functions}, in \emph{Frontiers in Number Theory, Physics, and Geometry I}, Springer,
2006, 365--377.

\bibitem{PlancherelPolya} M.~Plancherel, G.~P\'olya, \emph{Fonctions enti\`eres et
int\'egrales de Fourier multiples}, Comment.\ Math.\ Helv.\ 9 (1937), 224--248; 10 (1938),
110--163.

\bibitem{Seip} K.~Seip, \emph{Interpolation and Sampling in Spaces of Analytic Functions},
Univ.\ Lecture Ser.\ 33, Amer.\ Math.\ Soc., 2004.

\bibitem{Selberg1942} A.~Selberg, \emph{On the zeros of Riemann's zeta-function}, Skr.\
Norske Vid.\ Akad.\ Oslo I (1942), no.~10, 1--59.

\bibitem{Selberg} A.~Selberg, \emph{On the remainder in the formula for $N(T)$, the number
of zeros of $\zeta(s)$ in the strip $0<t<T$}, Avh.\ Norske Vid.\ Akad.\ Oslo I (1944),
no.~1, 1--27.

\bibitem{Selberg1946} A.~Selberg, \emph{Contributions to the theory of the Riemann
zeta-function}, Arch.\ Math.\ Naturvid.\ 48 (1946), no.~5, 89--155.

\bibitem{Suzuki} M.~Suzuki, \emph{On the Hilbert space derived from the Weil
distribution}, Canad.\ J.\ Math., published online 2025,
doi:10.4153/S0008414X25101739; arXiv:2301.00421.

\bibitem{Titchmarsh} E.~C.~Titchmarsh, \emph{The Theory of the Riemann Zeta-Function},
2nd ed., revised by D.~R.~Heath-Brown, Oxford Univ.\ Press, 1986.

\bibitem{Weil1952} A.~Weil, \emph{Sur les ``formules explicites'' de la th\'eorie des
nombres premiers}, Comm.\ S\'em.\ Math.\ Univ.\ Lund [Medd.\ Lunds Univ.\ Mat.\ Sem.],
Tome Suppl\'ementaire (1952), 252--265.

\bibitem{Yoshida} H.~Yoshida, \emph{On Hermitian forms attached to zeta functions}, in
\emph{Zeta Functions in Geometry}, Adv.\ Stud.\ Pure Math.\ 21 (1992), 281--325.

\end{thebibliography}

\end{document}
